\subsection*{A simple CM abelian surface in an actual finite-cover Prym}

Work over $L=\mathbf Q(i,\zeta_5)$. Let
\[
 C_{20}:x^{20}+y^{20}=z^{20},\qquad
 a:C_{20}\longrightarrow C_4,\quad(x:y:z)\longmapsto(x^5:y^5:z^5).
\]
The target $C_4$ and the map $f_4=[4]q_x:C_4\to E_0$ are those of
F.3.16. The degree of $a$ is25; its deck group over $L$ is
$(\mu_5)^2$, and it ramifies only above the coordinate-zero points of
$C_4$. F.3.17 proves that all those points map under $q_x$ into $D[4]$.
Consequently
\[
 f_{20}=f_4\circ a:C_{20}\to E_0
\]
has degree800 and is branched only over $O$. Thus
\[
 Y_{20}=C_{20}\setminus f_{20}^{-1}(O)\longrightarrow
 U=E_0\setminus\{O\}
\]
is an actual finite etale cover. The smooth plane curve $C_{20}$ has
genus171.

Let $X_5$ be the smooth proper genus-two curve with affine equation
\[
 X_5:\quad Y^2=1-X^5,
\]
and let $\infty$ be its unique point at infinity. The formulas
\[
 q:C_{20}\longrightarrow X_5,\qquad
 X=(x/z)^4,\quad Y=(y/z)^{10}
\]
give a morphism of smooth proper curves. It has degree40. Indeed on
function fields $L(C_{20})=L(u,v)$ with $u^{20}+v^{20}=1$, while
$L(X_5)=L(u^4,v^{10})$. The extension has degree at most40 by the
displayed fourth and tenth roots. The40 distinct automorphisms
$(u,v)\mapsto(\alpha u,\beta v)$ for
$\alpha\in\mu_4$, $\beta\in\mu_{10}$ fix this subfield; hence the
degree is exactly40.

\paragraph{CM and absolute simplicity, with the algebra embedding obstruction proved.}
Put $B=J_{X_5}$. The automorphism
\[
 \tau:(X,Y)\longmapsto(\zeta_5X,Y)
\]
fixes $\infty$ and induces an endomorphism of $B$. Its eigenvalues on
the basis $dX/Y,X\,dX/Y$ of regular differentials are
$\zeta_5,\zeta_5^2$; on the complex-conjugate part they are
$\zeta_5^{-1},\zeta_5^{-2}$. Thus its minimal polynomial on rational
first homology is $\Phi_5$. The faithful homology representation of
endomorphisms gives an actual embedding
\[
 E=\mathbf Q(\zeta_5)\hookrightarrow\operatorname{End}^0_L(B).
\]
This has degree4, equal to $2\dim B$, so $B$ is of CM type over $L$.

We prove absolute simplicity without assuming an unproved description
of the entire endomorphism algebra. A nonsimple abelian surface over
$\overline L$ is isogenous to a product of two elliptic curves.
If the factors are nonisogenous, its endomorphism algebra is
$D_1\times D_2$, where each $D_i$ is $\mathbf Q$ or an imaginary
quadratic field. A unital embedding of the degree-four field $E$ into
this product would give unital embeddings into each factor, which is
impossible by their dimensions.

For two isogenous elliptic factors, the algebra is $M_2(D)$ with
$D=\mathbf Q$ or an imaginary quadratic field. When $D=\mathbf Q$,
every matrix satisfies its degree-two characteristic polynomial, so
the image of a primitive element of $E$ cannot have degree four.
When $D$ is imaginary quadratic, its central copy in $M_2(D)$ commutes
with $E$. The field $E$ is cyclic quartic and its only quadratic
subfield is the real field $\mathbf Q(\sqrt5)$: the Galois group
$(\mathbf Z/5)^\times$ is cyclic of order4 and complex conjugation
is its unique subgroup of order2. Hence $D$ is not contained in $E$,
and $E\cap D=\mathbf Q$. Since $D/\mathbf Q$ is quadratic,
$E\otimes_{\mathbf Q}D$ is a field of degree8 over $\mathbf Q$.
Multiplication of the commuting copies defines a unital homomorphism
\[
 E\otimes_{\mathbf Q}D\longrightarrow M_2(D).
\]
Its source is a field, so this homomorphism is injective. Both sides
have dimension8 over $\mathbf Q$, making it an isomorphism of
algebras. That is impossible: the source is commutative and the
target contains the explicitly noncommuting matrix units
$e_{12},e_{21}$, whose two products are $e_{11},e_{22}$.
This rules out every elliptic isogeny decomposition and proves that
$B$ is absolutely simple.

\paragraph{The actual Prym inclusion.}
The norm and pullback maps satisfy
\[
 q_*q^*=[40]:B\to B.
\]
Every homomorphism $B\to E_0$ is zero: a nonzero such homomorphism
would be surjective, and its connected kernel would be a
one-dimensional abelian subvariety of the absolutely simple surface
$B$. In particular the actual composite
\[
 B\xrightarrow{q^*}J_{C_{20}}\xrightarrow{f_{20,*}}E_0
\]
is zero. The image of $q^*$ is connected and therefore belongs to
$P=\ker(f_{20,*})^0$. The identity $q_*q^*=[40]$ shows that its kernel
is finite, so it is an abelian subvariety of $P$ isogenous to $B$.
For every prime $\ell$ it gives the exact injection
\[
 T_\ell B\xrightarrow{q^*}T_\ell P\subset T_\ell J_{C_{20}},
\]
including $\ell=2,5$: applying $q_*$ multiplies a kernel vector by40,
and Tate modules are torsion-free. Rationally, $\tfrac1{40}q_*$ is a
retraction onto this factor.

For any $\ell\equiv1\pmod5$, the algebra
$E\otimes\mathbf Q_\ell\simeq\mathbf Q_\ell^4$ acts on
$V_\ell B$ with four one-dimensional eigenspaces. Every endomorphism
of this $E$-action is defined over $L$, so all four lines are
$G_L$-stable. Their images in $V_\ell P$ are genuine weight-one
CM lines in this finite-cover Prym.

The map to the punctured genus-two curve is also defined on the same
open cover: $q^{-1}(\infty)$ consists precisely of the coordinate-zero
points $z=0$, and these have already been removed in $Y_{20}$.
Thus $q$ restricts to
\[
 Y_{20}\longrightarrow X_5\setminus\{\infty\}.
\]
The explicit target has genus two, and its Jacobian has just been
proved to have no elliptic isogeny factor. This identifies the precise
geometric target beyond the elliptic maps in F.3.16. The construction
provides no procyclic lift of these Tate-module lines with open
arithmetic normalizer. LM6 A.22 remains unresolved.
